\subsection*{Why you cannot "just overlap" without events}

The \texttt{swap} pattern works because one thread group controls both loading and computing --- it knows the order. Role specialization (Section~\ref{sec:3-8}) puts them in \textbf{different} regions with different program counters. They cannot share a \texttt{swap} schedule; they need a signaling mechanism.

The picture below contrasts a strict load-then-compute staircase with double-buffered overlap: same logical work, less idle time.

Figure~\ref{fig:f08-fig1_pipeline_timeline} puts both schedules on a timeline.

\begin{figure}[htbp]\centering
\includegraphics[width=\linewidth,height=0.56\textheight,keepaspectratio]{f08-fig1_pipeline_timeline_light.png}
\caption{The same four K tiles run sequentially or with two-buffer overlap. Timings are schematic and include pipeline startup and drain.}
\label{fig:f08-fig1_pipeline_timeline}
\end{figure}
